% ===================================================================
%  4. ANALISI DEI REQUISITI E TRACCIABILITA'
% ===================================================================
\newpage
\section{Analisi dei Requisiti}

\subsection{Introduzione e Classificazione}
In questa sezione vengono elencati, classificati e descritti formalmente tutti i requisiti individuati durante la fase di analisi del dominio, dallo studio del capitolato d'appalto e dalle decisioni architetturali interne. I requisiti rappresentano i vincoli operativi e strutturali che il sistema \textbf{SmartOrder} deve soddisfare.

I requisiti sono identificati da un codice univoco secondo la seguente struttura:
\begin{center}
    \textbf{[TIPO]-[PRIORITÀ]\_[ID]}
\end{center}

\textbf{TIPO di requisito:}
\begin{itemize}
    \item \textbf{RF}: Requisito Funzionale (Cosa fa il sistema).
    \item \textbf{RQA}: Requisito Qualitativo / Non Funzionale (Come si comporta il sistema).
    \item \textbf{RP}: Requisito Prestazionale (Limiti e tempi di risposta).
    \item \textbf{RV}: Requisito di Vincolo (Scelte tecnologiche e architetturali obbligate).
\end{itemize}

\textbf{PRIORITÀ del requisito:}
\begin{itemize}
    \item \textbf{OB (Obbligatorio)}: Requisito fondamentale per il funzionamento base.
    \item \textbf{DE (Desiderabile)}: Requisito a valore aggiunto, utile ma non strettamente vitale.
    \item \textbf{OP (Opzionale)}: Possibile sviluppo futuro, non previsto in questa fase.
\end{itemize}

% -------------------------------------------------------------------
\newpage
\subsection{Requisiti Funzionali}

\subsubsection{Requisiti Funzionali obbligatori}
\renewcommand{\arraystretch}{1.15}
\begin{longtable}{|>{\raggedright}p{0.18\textwidth}|>{\raggedright\arraybackslash}p{0.5\textwidth}|>{\raggedright\arraybackslash}p{0.2\textwidth}|}
    \hline
    \rowcolor[gray]{0.9}
    \textbf{Codice} & \textbf{Descrizione} & \textbf{Fonti} \\
    \hline
    \endfirsthead
    \hline
    \rowcolor[gray]{0.9}
    \textbf{Codice} & \textbf{Descrizione} & \textbf{Fonti} \\
    \hline
    \endhead
    \hline
    RF-OB\_01 \hypertarget{rf-ob_1} & Il sistema deve permettere all'Utente di inserire il proprio indirizzo email nel campo dedicato per effettuare il login. & \hyperref[uc_01.1]{UC\_01.1} \\
    \hline
    RF-OB\_02 \hypertarget{rf-ob_2} & Il sistema deve permettere all'Utente di inserire la propria password in modo oscurato tramite il campo dedicato. & \hyperref[uc_01.2]{UC\_01.2} \\
    \hline
    RF-OB\_03 \hypertarget{rf-ob_3} & Il sistema deve fornire un pulsante o comando per confermare l'accesso e inviare le credenziali per la validazione. & \hyperref[uc_01]{UC\_01} \\
    \hline
    RF-OB\_04 \hypertarget{rf-ob_4} & Il sistema deve mostrare un errore se l'email inserita non ha un formato valido. & \hyperref[uc_01]{UC\_01} \\
    \hline
    RF-OB\_05 \hypertarget{rf-ob_5} & Il sistema deve mostrare un errore se l'email inserita non è presente nel database locale. & \hyperref[uc_01]{UC\_01} \\
    \hline
    RF-OB\_06 \hypertarget{rf-ob_6} & Il sistema deve mostrare un errore se la combinazione di email e password è sbagliata. & \hyperref[uc_01]{UC\_01} \\
    \hline
    RF-OB\_07 \hypertarget{rf-ob_7} & Dopo un login valido, il sistema deve creare la sessione e far accedere l'utente alla propria area (Cliente o Operatore). & \hyperref[uc_01]{UC\_01} \\
    \hline
    RF-OB\_08 \hypertarget{rf-ob_8} & Il sistema deve fornire un pulsante per effettuare il logout, invalidando e chiudendo la sessione corrente lato server. & \hyperref[uc_03]{UC\_03} \\
    \hline
    RF-OB\_09 \hypertarget{rf-ob_9} & Il sistema deve mostrare i dati aziendali del Cliente (ID e Ragione Sociale) in sola lettura. & \hyperref[uc_00.2]{UC\_00.2} \\
    \hline
    RF-OB\_10 \hypertarget{rf-ob_10} & Il sistema deve rendere accessibile la pagina Area Personale tramite un elemento di navigazione dedicato, consentendo all'Utente di consultare le informazioni del proprio profilo in base al ruolo assegnato. & \hyperref[uc_04]{UC\_04} \\
    \hline
    RF-OB\_11 \hypertarget{rf-ob_11} & Il sistema deve mostrare all'Utente il proprio indirizzo email in sola lettura. & \hyperref[uc_04.1]{UC\_04.1} \\
    \hline
    RF-OB\_12 \hypertarget{rf-ob_12} & Il sistema deve esporre un comando di contatto che, alla sua attivazione, invochi il protocollo URI \textit{mailto:} per aprire il client di posta elettronica nativo dell'Utente con il campo destinatario pre-compilato con l'email di supporto aziendale. & \hyperref[uc_06]{UC\_06} \\
    \hline
    RF-OB\_13 \hypertarget{rf-ob_13} & Il sistema deve mostrare in tempo reale tutti i messaggi inviati dall'Utente nella chat corrente. & \hyperref[uc_07]{UC\_07} \\
    \hline
    RF-OB\_14 \hypertarget{rf-ob_14} & Il sistema deve mostrare in tempo reale le risposte testuali scritte dall'AI o dall'Operatore di supporto. & \hyperref[uc_07]{UC\_07} \\
    \hline
    RF-OB\_15 \hypertarget{rf-ob_15} & Il sistema deve fornire un campo di testo per l'inserimento e l'invio di messaggi testuali nella chat. & \hyperref[uc_08.1]{UC\_08.1} \\
    \hline
    RF-OB\_16 \hypertarget{rf-ob_16} & Il sistema deve bloccare l'invio e mostrare un errore se il messaggio di testo supera il limite di caratteri. & \hyperref[uc_08.1]{UC\_08.1} \\
    \hline
    RF-OB\_17 \hypertarget{rf-ob_17} & Il sistema deve permettere all'Utente la registrazione e l'invio di messaggi vocali tramite il microfono del dispositivo. & \hyperref[uc_08.2]{UC\_08.2} \\
    \hline
    RF-OB\_18 \hypertarget{rf-ob_18} & Il sistema deve convertire l'audio in testo usando un modulo Speech-to-Text. & \hyperref[uc_08.2]{UC\_08.2} \\
    \hline
    RF-OB\_19 \hypertarget{rf-ob_19} & Il sistema deve bloccare l'invio e mostrare un errore se il file audio pesa più di 10MB. & \hyperref[uc_08.2]{UC\_08.2}, Capitolato \\
    \hline
    RF-OB\_20 \hypertarget{rf-ob_20} & Il sistema deve richiedere esplicitamente i permessi hardware necessari (microfono, fotocamera) al momento del primo utilizzo della relativa funzionalità e, in caso di rifiuto da parte dell'Utente, mostrare un messaggio informativo che spieghi l'impossibilità di procedere senza tali permessi. & Interno \\
    \hline
    RF-OB\_21 \hypertarget{rf-ob_21} & Il Cliente deve poter resettare la chat con l'AI e svuotare il carrello attuale. & \hyperref[uc_10]{UC\_10} \\
    \hline
    RF-OB\_22 \hypertarget{rf-ob_22} & Il sistema deve disabilitare il comando di reset della sessione qualora sia presente un ticket di assistenza in stato ``Aperto'' o ``In lavorazione'' per la sessione corrente, impedendo all'Utente di azzerare il contesto durante un'assistenza attiva. & \hyperref[uc_10]{UC\_10}, Interno \\
    \hline
    RF-OB\_23 \hypertarget{rf-ob_23} & L'AI (LLM) deve capire quali prodotti vuole l'utente e in che quantità leggendo/ascoltando il messaggio. & Capitolato \\
    \hline
    RF-OB\_24 \hypertarget{rf-ob_24} & Il sistema deve cercare i prodotti nel database usando la ricerca vettoriale (Vector Search). & Interno \\
    \hline
    RF-OB\_25 \hypertarget{rf-ob_25} & Il sistema deve aggiungere in automatico i prodotti trovati al carrello se il Confidence Score calcolato dall'AI supera la soglia definita in \hyperlink{rp-ob_3}{RP-OB\_03}. & Interno \\
    \hline
    RF-OB\_26 \hypertarget{rf-ob_26} & Vincolo Assortimento: Il sistema deve ignorare i prodotti che non fanno parte del catalogo assegnato a quel Cliente (tabella \texttt{asscli}). & Interno \\
    \hline
    RF-OB\_27 \hypertarget{rf-ob_27} & Vincolo Stato: Il sistema deve impedire l'aggiunta di prodotti segnati come ``non disponibili'' o ``sospesi'' nel database (\texttt{anaart}). & Interno \\
    \hline
    RF-OB\_28 \hypertarget{rf-ob_28} & Deduzione Formati: L'AI deve convertire automaticamente le quantità richieste nelle giuste unità di vendita (es. da bottiglie a casse). & Interno \\
    \hline
    RF-OB\_29 \hypertarget{rf-ob_29} & Il sistema deve fornire una barra di ricerca per cercare manualmente i prodotti nel catalogo tramite input testuale. & \hyperref[uc_11]{UC\_11} \\
    \hline
    RF-OB\_30 \hypertarget{rf-ob_30} & Il sistema deve mostrare i risultati della ricerca in tempo reale (Live Search). & \hyperref[uc_11]{UC\_11} \\
    \hline
    RF-OB\_31 \hypertarget{rf-ob_31} & Il sistema deve mostrare per ogni prodotto trovato: Nome, ID, Linea, Famiglia e Modalità di Vendita. & \hyperref[uc_00.1]{UC\_00.1} \\
    \hline
    RF-OB\_32 \hypertarget{rf-ob_32} & L'Utente deve poter cliccare su un risultato per decidere quanti pezzi acquistarne. & \hyperref[uc_12]{UC\_12} \\
    \hline
    RF-OB\_33 \hypertarget{rf-ob_33} & L'Utente deve poter modificare la quantità usando i tasti +/- o scrivendo direttamente il numero. & \hyperref[uc_12]{UC\_12} \\
    \hline
    RF-OB\_34 \hypertarget{rf-ob_34} & Il sistema deve correggere in automatico la quantità se l'utente inserisce numeri negativi, zero o troppo grandi. & \hyperref[uc_12]{UC\_12} \\
    \hline
    RF-OB\_35 \hypertarget{rf-ob_35} & L'Utente deve poter confermare l'aggiunta del prodotto al carrello della sessione corrente. & \hyperref[uc_12]{UC\_12} \\
    \hline
    RF-OB\_36 \hypertarget{rf-ob_36} & Il sistema deve mostrare un avviso di ``Carrello Vuoto'' se non ci sono prodotti inseriti. & \hyperref[uc_13]{UC\_13} \\
    \hline
    RF-OB\_37 \hypertarget{rf-ob_37} & Il sistema deve mostrare i dettagli logistici e le quantità di tutti i prodotti presenti nel carrello. & \hyperref[uc_13]{UC\_13} \\
    \hline
    RF-OB\_38 \hypertarget{rf-ob_38} & L'Utente deve poter sbloccare la riga di un prodotto nel carrello per poterne modificare la quantità. & \hyperref[uc_14.2]{UC\_14.2} \\
    \hline
    RF-OB\_39 \hypertarget{rf-ob_39} & L'Utente deve poter salvare la nuova quantità modificata per un prodotto nel carrello. & \hyperref[uc_14]{UC\_14} \\
    \hline
    RF-OB\_40 \hypertarget{rf-ob_40} & L'Utente deve poter rimuovere definitivamente un articolo dal carrello. & \hyperref[uc_14.1]{UC\_14.1} \\
    \hline
    RF-OB\_41 \hypertarget{rf-ob_41} & Il sistema deve disabilitare il pulsante per inviare l'ordine se il carrello è vuoto. & \hyperref[uc_15]{UC\_15} \\
    \hline
    RF-OB\_42 \hypertarget{rf-ob_42} & L'Utente (Cliente o Operatore in subentro) deve poter chiudere il carrello e inviare l'ordine definitivo. & \hyperref[uc_15]{UC\_15} \\
    \hline
    RF-OB\_43 \hypertarget{rf-ob_43} & Il sistema deve bloccare l'invio dell'ordine e segnalare gli errori se ci sono problemi con le regole del gestionale (es. prodotto fuori catalogo). & \hyperref[uc_15]{UC\_15} \\
    \hline
    RF-OB\_44 \hypertarget{rf-ob_44} & Dopo l'invio dell'ordine, il sistema deve svuotare la sessione e reindirizzare l'Utente alla vista appropriata al proprio ruolo: il Cliente allo Storico Ordini, l'Operatore alla chiusura del ticket e dismissione della console di subentro. & \hyperref[uc_15]{UC\_15}, Interno \\
    \hline
    RF-OB\_45 \hypertarget{rf-ob_45} & Il sistema deve creare un file JSON formattato contenente tutti i dati dell'ordine appena chiuso. & \hyperref[uc_15]{UC\_15}, Capitolato \\
    \hline
    RF-OB\_46 \hypertarget{rf-ob_46} & Il Cliente deve poter visualizzare la lista dei propri ordini passati, ordinati per data. & \hyperref[uc_16]{UC\_16} \\
    \hline
    RF-OB\_47 \hypertarget{rf-ob_47} & Il sistema deve mostrare per ogni ordine storico: ID, Data, totale degli articoli e una piccola anteprima dei prodotti. & \hyperref[uc_00.3]{UC\_00.3}, \hyperref[uc_16]{UC\_16} \\
    \hline
    RF-OB\_48 \hypertarget{rf-ob_48} & Il Cliente deve poter cliccare su un ordine passato per vederne tutti i dettagli logistici. & \hyperref[uc_18]{UC\_18} \\
    \hline
    RF-OB\_49 \hypertarget{rf-ob_49} & Il sistema deve mostrare in sola lettura la chat originale tra utente e AI che ha generato quell'ordine specifico. & \hyperref[uc_00.4]{UC\_00.4} \\
    \hline
    RF-OB\_50 \hypertarget{rf-ob_50} & Il Cliente deve poter richiedere l'aiuto di un operatore umano, mettendo in pausa l'AI. & \hyperref[uc_19]{UC\_19} \\
    \hline
    RF-OB\_51 \hypertarget{rf-ob_51} & L'Operatore deve poter accedere alla dashboard con la lista dei Ticket Aperti (richieste d'aiuto). & \hyperref[uc_21]{UC\_21} \\
    \hline
    RF-OB\_52 \hypertarget{rf-ob_52} & Nella lista dei ticket, il sistema deve mostrare da quanto tempo il Cliente sta aspettando una risposta. & \hyperref[uc_21.5]{UC\_21.5} \\
    \hline
    RF-OB\_53 \hypertarget{rf-ob_53} & L'Operatore deve poter prendere in carico un ticket bloccandolo per gli altri colleghi (Lock) per aprire la console di assistenza. & \hyperref[uc_22]{UC\_22} \\
    \hline
    RF-OB\_54 \hypertarget{rf-ob_54} & L'Operatore deve poter scrivere e inviare messaggi in chat direttamente al Cliente per fornirgli assistenza. & \hyperref[uc_08.1]{UC\_08.1} \\
    \hline
    RF-OB\_55 \hypertarget{rf-ob_55} & L'Operatore deve poter usare la ricerca per aggiungere forzatamente prodotti al carrello del Cliente che sta aiutando. & \hyperref[uc_11]{UC\_11}, \hyperref[uc_12]{UC\_12} \\
    \hline
    RF-OB\_56 \hypertarget{rf-ob_56} & L'Operatore deve poter modificare le quantità o rimuovere prodotti dal carrello del Cliente. & \hyperref[uc_13]{UC\_13}, \hyperref[uc_14]{UC\_14} \\
    \hline
    RF-OB\_57 \hypertarget{rf-ob_57} & L'Operatore deve poter chiudere il ticket per segnare il problema come risolto e sbloccare la sessione (rimuovere il Lock). & \hyperref[uc_22.2]{UC\_22.2} \\
    \hline
    RF-OB\_58 \hypertarget{rf-ob_58} & L'Utente (Cliente o Operatore) deve poter chiudere il ticket da solo se risolve il problema o non ha più bisogno di aiuto. & \hyperref[uc_20]{UC\_20} \\
    \hline
    RF-OB\_59 \hypertarget{rf-ob_59} & L'Operatore deve poter accedere alla dashboard che mostra tutti gli ordini inviati da tutti i Clienti. & \hyperref[uc_23]{UC\_23} \\
    \hline
    RF-OB\_60 \hypertarget{rf-ob_60} & L'Operatore deve poter cliccare su un ordine di qualsiasi cliente per aprirlo e vederne i dettagli. & \hyperref[uc_24]{UC\_24} \\
    \hline
    RF-OB\_61 \hypertarget{rf-ob_61} & L'Operatore deve poter cercare del testo solo all'interno di una colonna specifica (es. cercare solo l'ID Ordine). & \hyperref[uc_31]{UC\_31} \\
    \hline
    RF-OB\_62 \hypertarget{rf-ob_62} & L'Operatore deve poter ordinare le tabelle in modo crescente o decrescente cliccando sui nomi delle colonne. & \hyperref[uc_29]{UC\_29} \\
    \hline
    RF-OB\_63 \hypertarget{rf-ob_63} & L'Operatore deve poter filtrare i dati delle tabelle scegliendo una data di inizio e di fine. & \hyperref[uc_30]{UC\_30} \\
    \hline
    RF-OB\_64 \hypertarget{rf-ob_64} & Il sistema deve salvare nel database tutto lo storico delle chat e delle decisioni prese dall'AI. & Interno, Capitolato \\
    \hline
    RF-OB\_65 \hypertarget{rf-ob_65} & Il sistema deve mostrare la percentuale di sicurezza dell'AI (Confidence Score) solo nelle schermate dedicate all'Operatore. & \hyperref[uc_22.1]{UC\_22.1}, \hyperref[uc_24.1]{UC\_24.1} \\
    \hline
    RF-OB\_66 \hypertarget{rf-ob_66} & Il sistema deve mostrare schermate di caricamento (spinner o barre di avanzamento) durante le operazioni che richiedono un tempo di elaborazione superiore a 1 secondo. & Interno \\
    \hline
    RF-OB\_67 \hypertarget{rf-ob_67} & Le API del server (Backend) devono rifiutare qualsiasi richiesta priva di un token di accesso valido per la sessione corrente. & Interno \\
    \hline
    RF-OB\_68 \hypertarget{rf-ob_68} & Il sistema deve caricare i dati delle tabelle in modo incrementale tramite paginazione con pagine da massimo 50 record, per evitare il blocco dell'interfaccia in presenza di grandi volumi di dati. & Interno \\
    \hline
    \multicolumn{3}{c}{} \\
    \caption{Requisiti Funzionali Obbligatori}\hypertarget{tab:funzionali_obbligatori}\\
\end{longtable}

\subsubsection{Requisiti Funzionali desiderabili}
\renewcommand{\arraystretch}{1.15}
\begin{longtable}{|>{\raggedright}p{0.18\textwidth}|>{\raggedright\arraybackslash}p{0.5\textwidth}|>{\raggedright\arraybackslash}p{0.2\textwidth}|}
    \hline
    \rowcolor[gray]{0.9}
    \textbf{Codice} & \textbf{Descrizione} & \textbf{Fonti} \\
    \hline
    \endfirsthead
    \hline
    \rowcolor[gray]{0.9}
    \textbf{Codice} & \textbf{Descrizione} & \textbf{Fonti} \\
    \hline
    \endhead
    \hline
    RF-DE\_01 \hypertarget{rf-de_1} & Il sistema deve fornire un link per accedere alla pagina di recupero della password dimenticata. & \hyperref[uc_02]{UC\_02} \\
    \hline
    RF-DE\_02 \hypertarget{rf-de_2} & L'Utente deve poter richiedere il reset della password inserendo la propria email. & \hyperref[uc_02]{UC\_02}, \hyperref[uc_02.1]{UC\_02.1} \\
    \hline
    RF-DE\_03 \hypertarget{rf-de_3} & Il sistema deve permettere all'Utente loggato di avviare la procedura per cambiare la propria password. & \hyperref[uc_04.2]{UC\_04.2} \\
    \hline
    RF-DE\_04 \hypertarget{rf-de_4} & L'Utente deve poter inviare la vecchia password e la nuova password scelta per salvarla nel sistema. & \hyperref[uc_05.1]{UC\_05.1}, \hyperref[uc_05.2]{UC\_05.2}, \hyperref[uc_05.3]{UC\_05.3} \\
    \hline
    RF-DE\_05 \hypertarget{rf-de_5} & Il sistema deve bloccare il cambio password se la vecchia password è errata o se le due nuove password inserite non coincidono. & \hyperref[uc_05.1]{UC\_05.1}, \hyperref[uc_05.3]{UC\_05.3} \\
    \hline
    RF-DE\_06 \hypertarget{rf-de_6} & Il sistema deve permettere l'invio di file immagine nella chat, mediante scatto fotografico o selezione dalla galleria del dispositivo. & \hyperref[uc_08.3]{UC\_08.3} \\
    \hline
    RF-DE\_07 \hypertarget{rf-de_7} & Il sistema deve bloccare l'invio di immagini e mostrare un errore se il file pesa più di 5MB. & \hyperref[uc_08.3]{UC\_08.3}, Interno \\
    \hline
    RF-DE\_08 \hypertarget{rf-de_8} & Il Cliente deve poter valutare la qualità delle risposte dell'AI lasciando un feedback (es. pollice su / pollice giù). & \hyperref[uc_09.1]{UC\_09.1}, \hyperref[uc_09.2]{UC\_09.2} \\
    \hline
    RF-DE\_09 \hypertarget{rf-de_9} & Se il Cliente dà un voto negativo all'AI, il sistema deve richiedere che indichi il tipo di errore (tramite tag predisposti). & \hyperref[uc_09.2]{UC\_09.2} \\
    \hline
    RF-DE\_10 \hypertarget{rf-de_10} & Il Cliente deve poter aggiungere un testo libero per spiegare meglio il motivo dell'errore segnalato. & \hyperref[uc_09.2]{UC\_09.2} \\
    \hline
    RF-DE\_11 \hypertarget{rf-de_11} & Il Cliente deve poter scaricare l'elenco dei propri ordini storici in un file formato CSV. & \hyperref[uc_17]{UC\_17} \\
    \hline
    RF-DE\_12 \hypertarget{rf-de_12} & L'Operatore deve poter accedere alla dashboard per il controllo degli errori e l'addestramento dell'AI (Ottimizzazione AI). & \hyperref[uc_25]{UC\_25} \\
    \hline
    RF-DE\_13 \hypertarget{rf-de_13} & Il sistema deve mostrare all'Operatore la lista degli errori segnalati dai clienti e le correzioni manuali fatte ai carrelli. & \hyperref[uc_25]{UC\_25} \\
    \hline
    RF-DE\_14 \hypertarget{rf-de_14} & L'Operatore deve poter aprire un singolo errore AI per rileggere tutta la chat originale tra il Cliente e il bot. & \hyperref[uc_26]{UC\_26} \\
    \hline
    RF-DE\_15 \hypertarget{rf-de_15} & L'Operatore deve poter approvare una segnalazione di errore, confermando che i dati verranno usati per riaddestrare l'AI. & \hyperref[uc_27.1]{UC\_27.1} \\
    \hline
    RF-DE\_16 \hypertarget{rf-de_16} & L'Operatore deve poter scartare una segnalazione se non si tratta di un vero errore dell'AI (es. l'utente ha scritto male). & \hyperref[uc_27.2]{UC\_27.2} \\
    \hline
    RF-DE\_17 \hypertarget{rf-de_17} & L'Operatore deve poter annullare l'approvazione di una segnalazione se l'aveva confermata per sbaglio (Rollback). & \hyperref[uc_28]{UC\_28} \\
    \hline
    \multicolumn{3}{c}{} \\
    \caption{Requisiti Funzionali Desiderabili}\hypertarget{tab:funzionali_desiderabili}\\
\end{longtable}

\subsubsection{Requisiti Funzionali opzionali}
\renewcommand{\arraystretch}{1.15}
\begin{longtable}{|>{\raggedright}p{0.18\textwidth}|>{\raggedright\arraybackslash}p{0.5\textwidth}|>{\raggedright\arraybackslash}p{0.2\textwidth}|}
    \hline
    \rowcolor[gray]{0.9}
    \textbf{Codice} & \textbf{Descrizione} & \textbf{Fonti} \\
    \hline
    \endfirsthead
    \hline
    \rowcolor[gray]{0.9}
    \textbf{Codice} & \textbf{Descrizione} & \textbf{Fonti} \\
    \hline
    \endhead
    \hline
    RF-OP\_01 \hypertarget{rf-op_1} & Il sistema deve permettere l'invio di messaggi d'ordine e la loro elaborazione tramite l'app di WhatsApp. & Interno \\
    \hline
    RF-OP\_02 \hypertarget{rf-op_2} & Il sistema deve permettere l'invio di messaggi d'ordine e la loro elaborazione tramite l'app di Telegram. & Interno \\
    \hline
    RF-OP\_03 \hypertarget{rf-op_3} & Il sistema deve generare grafici e statistiche generali sull'uso della piattaforma, mostrandoli all'Operatore. & Interno \\
    \hline
    \multicolumn{3}{c}{} \\
    \caption{Requisiti Funzionali Opzionali}\hypertarget{tab:funzionali_opzionali}\\
\end{longtable}

% -------------------------------------------------------------------
\newpage
\subsection{Requisiti Qualitativi (Non Funzionali)}
Questa sezione traccia i requisiti legati alle qualità del software, quali manutenibilità, readability e testabilità (identificati dal prefisso \textbf{RQA}).

\subsubsection{Requisiti Qualitativi obbligatori}
\renewcommand{\arraystretch}{1.15}
\begin{longtable}{|>{\raggedright}p{0.18\textwidth}|>{\raggedright\arraybackslash}p{0.5\textwidth}|>{\raggedright\arraybackslash}p{0.2\textwidth}|}
    \hline
    \rowcolor[gray]{0.9}
    \textbf{Codice} & \textbf{Descrizione} & \textbf{Fonti} \\
    \hline
    \endfirsthead
    \hline
    \rowcolor[gray]{0.9}
    \textbf{Codice} & \textbf{Descrizione} & \textbf{Fonti} \\
    \hline
    \endhead
    \hline
    RQA-OB\_01 \hypertarget{rqa-ob_1} & Il codice del sistema deve essere scritto seguendo gli standard di readability e manutenibilità definiti nelle Norme di Progetto. In particolare: funzioni non devono superare le 50 linee di codice, le costanti devono essere commentate, e i nomi di variabili e metodi devono essere significativi e seguire le convenzioni di naming. & Norme di Progetto, Interno \\
    \hline
    \multicolumn{3}{c}{} \\
    \caption{Requisiti Qualitativi obbligatori}\hypertarget{tab:qualitativi_obbligatori}\\
\end{longtable}

% -------------------------------------------------------------------
\newpage
\subsection{Requisiti Prestazionali}
Questa sezione raccoglie i limiti tecnici e i tempi di risposta che il sistema deve rispettare (identificati dal prefisso \textbf{RP}).

\subsubsection{Requisiti Prestazionali obbligatori}
\renewcommand{\arraystretch}{1.15}
\begin{longtable}{|>{\raggedright}p{0.18\textwidth}|>{\raggedright\arraybackslash}p{0.5\textwidth}|>{\raggedright\arraybackslash}p{0.2\textwidth}|}
    \hline
    \rowcolor[gray]{0.9}
    \textbf{Codice} & \textbf{Descrizione} & \textbf{Fonti} \\
    \hline
    \endfirsthead
    \hline
    \rowcolor[gray]{0.9}
    \textbf{Codice} & \textbf{Descrizione} & \textbf{Fonti} \\
    \hline
    \endhead
    \hline
    RP-OB\_01 \hypertarget{rp-ob_1} & Il modello AI (LLM) deve completare l'elaborazione del messaggio e rispondere entro massimo 20 secondi, misurati come valore mediano su un campione di almeno 50 richieste consecutive in assenza di traffico concorrente. & Interno, Capitolato \\
    \hline
    RP-OB\_02 \hypertarget{rp-ob_2} & Il sistema di trascrizione vocale (Speech-to-Text) deve riconoscere le parole con un Word Error Rate (WER) inferiore al 10\%, misurato su un campione di almeno 100 registrazioni audio in italiano standard, con qualità audio non inferiore a 16\,kHz mono e in assenza di rumore di fondo superiore a 40\,dB. & Interno \\
    \hline
    RP-OB\_03 \hypertarget{rp-ob_3} & La percentuale di sicurezza (Confidence Score) calcolata dall'AI deve superare il 70\% per poter aggiungere prodotti in automatico. & Interno \\
    \hline
    RP-OB\_04 \hypertarget{rp-ob_4} & L'applicazione web deve caricare gli elementi principali a schermo in meno di 3 secondi, misurati al 90° percentile su 10 caricamenti consecutivi da una connessione con banda minima di 20\,Mbit/s in download. & Interno \\
    \hline
    RP-OB\_05 \hypertarget{rp-ob_5} & Quando l'utente usa la barra di ricerca prodotti, i risultati devono comparire in meno di 500 millisecondi, misurati al 95° percentile su un database contenente almeno 2000 prodotti, con un massimo di 10 utenti concorrenti e connessione locale (LAN). & Interno \\
    \hline
    \multicolumn{3}{c}{} \\
    \caption{Requisiti Prestazionali obbligatori}\hypertarget{tab:prestazionali_obbligatori}\\
\end{longtable}

\subsubsection{Requisiti Prestazionali desiderabili}
\renewcommand{\arraystretch}{1.15}
\begin{longtable}{|>{\raggedright}p{0.18\textwidth}|>{\raggedright\arraybackslash}p{0.5\textwidth}|>{\raggedright\arraybackslash}p{0.2\textwidth}|}
    \hline
    \rowcolor[gray]{0.9}
    \textbf{Codice} & \textbf{Descrizione} & \textbf{Fonti} \\
    \hline
    \endfirsthead
    \hline
    \rowcolor[gray]{0.9}
    \textbf{Codice} & \textbf{Descrizione} & \textbf{Fonti} \\
    \hline
    \endhead
    \hline
    RP-DE\_01 \hypertarget{rp-de_1} & L'estrazione e il download in CSV dello storico ordini deve concludersi in massimo 5 secondi per liste con meno di 1000 ordini (dimensione stimata $\leq$ 2\,MB), misurato al 90° percentile su 10 esecuzioni consecutive in condizioni di rete locale (LAN). & Interno \\
    \hline
    \multicolumn{3}{c}{} \\
    \caption{Requisiti Prestazionali desiderabili}\hypertarget{tab:prestazionali_desiderabili}\\
\end{longtable}

% -------------------------------------------------------------------
\newpage
\subsection{Requisiti di Vincolo}
Questa sezione elenca gli obblighi e le scelte tecnologiche imposte dall'esterno a cui il sistema deve conformarsi (identificati dal prefisso \textbf{RV}).

\subsubsection{Requisiti di Vincolo obbligatori}
\renewcommand{\arraystretch}{1.15}
\begin{longtable}{|>{\raggedright}p{0.18\textwidth}|>{\raggedright\arraybackslash}p{0.5\textwidth}|>{\raggedright\arraybackslash}p{0.2\textwidth}|}
    \hline
    \rowcolor[gray]{0.9}
    \textbf{Codice} & \textbf{Descrizione} & \textbf{Fonti} \\
    \hline
    \endfirsthead
    \hline
    \rowcolor[gray]{0.9}
    \textbf{Codice} & \textbf{Descrizione} & \textbf{Fonti} \\
    \hline
    \endhead
    \hline
    RV-OB\_01 \hypertarget{rv-ob_1} & Il Frontend (la parte grafica usata dall'utente) e il Backend (il server) devono essere programmati come due sistemi separati che comunicano tra loro. & Interno, Norme di Progetto \\
    \hline
    RV-OB\_02 \hypertarget{rv-ob_2} & Il database della piattaforma web deve essere di tipo relazionale, nello specifico PostgreSQL. & Interno \\
    \hline
    RV-OB\_03 \hypertarget{rv-ob_3} & La comprensione del testo scritto dagli utenti deve essere gestita tramite tecnologie AI, nello specifico gli Embedding e la Vector Search. & Capitolato \\
    \hline
    RV-OB\_04 \hypertarget{rv-ob_4} & Il file JSON dell'ordine deve essere posizionato in una cartella o sistema condiviso in modo che il gestionale aziendale (ERP) possa leggerlo. & Interno (Logica ERP) \\
    \hline
    RV-OB\_05 \hypertarget{rv-ob_5} & Tutte le comunicazioni tra la webapp e il server devono usare esclusivamente connessioni sicure crittografate (HTTPS). & Interno \\
    \hline
    RV-OB\_06 \hypertarget{rv-ob_6} & Le password non devono mai essere salvate in chiaro, ma protette nel database usando l'algoritmo di hashing scrypt e salt. & Interno \\
    \hline
    RV-OB\_07 \hypertarget{rv-ob_7} & Il salvataggio nel database e l'uso delle email degli utenti deve rispettare le regole europee sulla privacy (GDPR). & Interno \\
    \hline
    RV-OB\_08 \hypertarget{rv-ob_8} & L'applicazione web deve funzionare correttamente sui seguenti browser desktop nelle versioni minime indicate: Google Chrome 120+, Mozilla Firefox 121+, Apple Safari 17+, Microsoft Edge 120+. & Interno \\
    \hline
    RV-OB\_09 \hypertarget{rv-ob_9} & L'applicazione web deve funzionare correttamente sui seguenti browser mobile nelle versioni minime indicate: Chrome per Android 120+, Safari per iOS 17+. & Interno \\
    \hline
    RV-OB\_10 \hypertarget{rv-ob_10} & L'interfaccia del Cliente deve adattarsi a schermi con larghezza minima di 375\,px, con approccio \textit{Mobile-first}: tutti i controlli interattivi devono avere un'area toccabile di almeno 44$\times$44\,px e l'interfaccia non deve presentare scroll orizzontale su alcun viewport supportato. & Interno \\
    \hline
    RV-OB\_11 \hypertarget{rv-ob_11} & L'interfaccia dell'Operatore è progettata per essere usata su monitor da computer con larghezza minima di 1024\,px, con approccio \textit{Desktop-first}: tutti i pannelli e i controlli devono essere completamente fruibili su questa larghezza senza degradazioni di layout. & Interno \\
    \hline
    \multicolumn{3}{c}{} \\
    \caption{Requisiti di Vincolo obbligatori}\hypertarget{tab:vincolo_obbligatori}\\
\end{longtable}
